\section{Background and definition}
\label{sec:background}

\subsection{What a harness is}
\label{subsec:definition}

An agent harness is the software layer between a language model and a task environment that
repeatedly assembles the model's input, executes the model's chosen actions, and decides whether to
continue. Let $M$ be a model, a map from input to output tokens together with its decoding
configuration, and let $E$ be a task environment: a state space, executable operations with observable
results, and a task specification. $M$ cannot act on $E$ and $E$ cannot read $M$; the harness $H$
connects them by implementing eight functions. It (i) assembles input, (ii) exposes and executes
actions, (iii) decides continuation and termination, (iv) holds state beyond one model call,
(v) verifies and repairs, (vi) enforces budgets and permissions, (vii) isolates execution, and
(viii) records and governs what happened. An agent is the run-time triple $(M, H, E)$, which makes the
harness effect well posed: hold $M$ and $E$ fixed and vary $H$.

Clauses (i) to (iii) are necessary; without them a layer is a prompt, a tool server, or a pipeline.
Clauses (iv) to (viii) may be trivial, and a trivial value is still a value: a harness that runs the
model in the host process is coded \texttt{none} on execution isolation, not \texttt{not\_reported}.
The minimal harness formats the transcript, executes one parsed action, and stops on a finish action
or a step cap; ReAct \citep{yao2023react} is exactly that, and is the earliest system in the corpus.
Later systems elaborate particular clauses. Reflexion keeps verbal reflections on task feedback in an
episodic memory \citep{shinn2023reflexion}, and CodeAct makes executable Python code the action space
\citep{wang2024codeact}. AutoGen and MetaGPT run several agents that converse or hold assigned roles
\citep{wu2023autogen,hong2023metagpt}, and Agent S and OS-Copilot act through a computer's interface,
evaluated on OSWorld and GAIA \citep{agashe2024agents,wu2024oscopilot,xie2024osworld,mialon2023gaia}.
The unit of analysis is the system, not the paper: coding papers would count a popular harness once
per publication and a harness with no paper not at all.

One test decides every boundary case: whether changing the thing in question, with everything else
fixed, would change what the model sees, what it can do, or when it stops. Table~\ref{tab:boundary} applies
it where the definition is most often probed. The pattern is that a part is not a system, while the
loop a framework or benchmark ships is one.

\begin{table}[t]
  \caption{Boundary cases and how the definition decides them.}
  \label{tab:boundary}
  \footnotesize
  \begin{tabular}{@{}p{0.40\linewidth}p{0.29\linewidth}p{0.25\linewidth}@{}}
    \toprule
    Case & Decision & Reason \\
    \midrule
    Weights, decoding settings, trained call syntax & Model, not harness & Fixed inside a forward
    pass \\
    Tool schemas, output parser, re-prompts & Harness & Change what the model sees \\
    Reasoning model that stops thinking itself & Harness still owns (iii) & Harness decides on
    another call \\
    \midrule
    Prompt; chain-of-thought with self-consistency; tree of thoughts \citep{yao2023tot} & Not a system & No action on an environment \\
    ReAct \citep{yao2023react} & System & Parse, execute, append, repeat \\
    SWE-agent with GPT-4 on SWE-bench Lite & An agent; SWE-agent is the system
    \citep{yang2024sweagent} & An agent is $(M, H, E)$ \\
    \midrule
    Framework library, e.g.\ LangGraph \citep{langgraph} or CrewAI \citep{crewai} & Coded via its default agent &
    A kit, not a loop \\
    Container the harness launches \citep{wang2024openhands} & Part of the harness & The harness
    configures it \\
    Sandbox product, protocol server, vector store, tracer & Not a system & One clause, not the
    loop \\
    \midrule
    Non-looping benchmark baseline \citep{jimenez2024swebench} & Not a system & It is $E$ \\
    Looping benchmark baseline \citep{zhou2023webarena,xie2024osworld,yao2024taubench,chezelles2024browsergym} & System & Loops a model over $E$ \\
    Tool-use fine-tuning \citep{schick2023toolformer} & Out of scope unless it adds a loop &
    Changes $M$, not $H$ \\
    Closed vendor product \citep{claudecode} & System if codable from documents & Defined by
    behavior \\
    \bottomrule
  \end{tabular}
\end{table}

\subsection{The nine layers}
\label{subsec:layers}

The eight clauses map onto coding layers A to H, plus a metadata layer M, for 38 dimensions.
Context assembly (A, four dimensions) governs what the model sees at each step, and tool interface
(B, five) what it may do and in what syntax. Control loop (C, five) governs what happens after a
step, including topology, delegation, and human checkpoints. Memory and state (D, three) covers what
persists within and across runs, verification and repair (E, three) whether an oracle checks the work,
and budget and termination (F, three) step, cost, and time bounds. Sandbox and environment (G, four)
covers isolation, filesystem, network, and permissions; observability and governance (H, four) covers
tracing, replay, evaluation hooks, and guardrails. Layer M (seven) describes the system record, such
as domain, openness, and primary artifact, and is excluded wherever a claim is about design. Three
placements are deliberate. Termination sits in F rather than C so that termination and budget are
counted together. Filesystem access and network policy sit apart from isolation because released
systems configure them independently. Observability and governance share H because in this corpus
they are reported together or not at all.

\subsection{The Three-State Cell Contract}
\label{subsec:cellstates}

The paper's measurement vocabulary is the Three-State Cell Contract. A cell is one system against one
dimension and holds exactly one state. A \emph{value} is backed by a verbatim quote and a re-openable
locator, and includes an absence found where the feature would be declared (50.8\% of the 47,728
cells). \texttt{not\_reported} means the sources were read and are silent (48.9\%), and
\texttt{unresolved} means the cell could not be settled (0.3\%). Silence is data about documents,
never absence in a system: the sources are silent on network policy for an estimated 93.3\% of
systems, which does not mean that 93.3\% of systems lack one. An \texttt{unresolved} cell is an
admission about our coding, excluded from every denominator and never pooled with
\texttt{not\_reported}, because pooling would inflate the silence rate by exactly what we failed to
code. Every share is labeled weighted, describing the frame, or unweighted, describing the coded set,
and carries its denominator. Section~\ref{sec:methods} gives the coding rule that enforces the
contract.
